\begin{tikzpicture}[font=\scriptsize, >=latex]
% laundry timeline: single-cycle vs pipelined
\node[font=\footnotesize\bfseries] at (6,5.6)
  {same 5 instructions, two machines};
% single cycle row
\node[left, font=\tiny] at (0.4,4.9) {single-cycle};
\foreach \i in {0,...,4}{
  \draw[thick, fill=accentsoft] (0.6+\i*1.0,4.5) rectangle (1.55+\i*1.0,5.15);}
\node[font=\tiny] at (3.1,4.25) {each box = one slow clock};
% pipeline lanes
\node[left, font=\tiny] at (0.4,2.6) {pipelined};
\foreach \lane/\lbl in {0/IF,1/ID,2/EX,3/MEM,4/WB}{
  \node[left, font=\tiny] at (6.9,3.4-\lane*0.5) {\lbl};}
\foreach \i in {0,...,4}{
  \foreach \lane in {0,...,4}{
    \pgfmathtruncatemacro{\s}{\i+\lane}
    \draw[thick, fill=accentsoft] (7.0+\s*0.45,3.4-\lane*0.5)
      rectangle (7.45+\s*0.45,3.9-\lane*0.5);}}
\node[font=\tiny] at (8.9,0.4) {instr $n$ in lane $n$ --- five at once};
\node[align=left, text=slate, text width=11cm] at (6.5,-0.7)
  {Left: 5 instructions in 5 long periods. Right: the same five
   overlap --- the pipeline finishes one instruction per short
   period once full. Latency per instruction is unchanged; the
   throughput is what multiplies.};
\end{tikzpicture}
